\section{Kernels and the geometry of available slots}\label{sec:geometry}

The fresh walk moves through a graph whose available slots change with
physical time.  We need two deterministic ways to control it: a Nash
inequality converts dissipation into diffusion, and flows to infinity
bound the pairing of a source mass with a test function by its Dirichlet
energy.  The first source bound permits many sparse starting points;
the second becomes sharper when a single source has already diffused.
We obtain both bounds on a fixed collection of nearest-neighbor
sublattices contained in the interaction graph.  The walk itself always
uses all interaction bonds.

\subsection{Coset coordinates and the killed kernels}

Choose linearly independent vectors \(a_1,a_2,a_3\) with \(J(a_j)>0\),
and let \(A\) be the integer matrix with these columns.  Such a choice
exists because the support of \(J\) generates \(\Z^3\).
The sublattice \(A\Z^3\) has finite index \(|\det A|\).  Fix
representatives \((u_\gamma)_{\gamma\in\mathfrak C}\) of its cosets and
write, uniquely,
\begin{equation}\label{eq:coset-coordinates}
 x=u_\gamma+A\bar x,
 \qquad \gamma\in\mathfrak C,\quad \bar x\in\Z^3.
\end{equation}
For \(z\in\Z^3\) and \(R\ge0\), put
\[
 B_R(z)=\{w\in\Z^3:|w-z|_\infty\le R\},
 \qquad
 \mathfrak B_R(z)=\{x\in\Z^3:\bar x\in B_R(z)\}.
\]
Thus \(\mathfrak B_R(z)\) contains a copy of the same coordinate ball
in every coset.  Each coset contains all nearest-neighbor edges of its
\(\bar x\)-coordinate lattice, since these edges are the displacements
\(\pm a_j\).  There is also a fixed integer \(B_*\ge1\) such that
\begin{equation}\label{eq:bond-coordinate-range}
 J(y-x)>0\quad\Longrightarrow\quad
 |\bar y-\bar x|_\infty\le B_*.
\end{equation}
Indeed, \(\bar y-\bar x=A^{-1}(y-x+u_\gamma-u_{\gamma'})\), and the
right side has only finitely many possibilities.  Constants below may
depend on this choice of coordinates, \(S,J\), and fixed geometric
radii, but never on the volume, pins, or subsequent scales.

Fix an exposure \(\mathbf a\) from Section~\ref{sec:pins} and a cut
\(\theta\).  Let \(U\) be the one-period fresh-walk kernel from the
holes at this cut back to the same holes, killed whenever it visits
\(\partial\Lambda\), including at its starting point.  The kernel
has no killing at other pins.  We use column-forward matrices: \(Uv\)
is the output mass when the input mass is \(v\).  Define
\begin{equation}\label{eq:one-period}
 U_+=U,\qquad U_-=U^*,\qquad
 Q=\frac{U+U^*}{2},\qquad A_Q=I-Q.
\end{equation}
The ordered kernel $U$ need not be self-adjoint. Its quadratic form
on real vectors agrees with that of $Q$, so the symmetric part is the
operator needed for the return-energy estimate.
The adjoint traversal reverses the order of all factors.  Between
prescribed changes the heat factors are symmetric; transports reverse
by their inverse bijections, and the adjoint of a uniform seam completion
is a uniform inverse completion.  Killing projections accompany these
factors in both directions.

Every interval kernel is an average of partial permutation matrices:
sample the simultaneous trajectories of all holes and retain precisely
those trajectories that survive the killing.  Consequently it is
nonnegative, with every row and column sum at most one, and contracts
counting \(1\)-, \(2\)-, and supremum norms.  For example, if \(K\)
is such a kernel, Cauchy--Schwarz gives
\[
 \sum_i|(Kv)_i|^2
 \le\sum_i\Bigl(\sum_jK_{ij}\Bigr)
              \Bigl(\sum_jK_{ij}|v_j|^2\Bigr)
 \le\sum_j|v_j|^2.
\]
These properties hold for adjoints, products, and further spatial
killing.  In particular \(Q\) is self-adjoint and a contraction, so
\(A_Q\) is positive semidefinite.

A split of a traversal always uses matching hole layers.  At a prescribed
transport or seam we distinguish its two sides and assign the operation
to exactly one of the adjoining factors.  Thus a full forward traversal
from the post-seam cut at zero includes its seam at the end.  Reversing
a traversal reverses this same ordered list.  Time integrals ignore the
finitely many instantaneous operations.

At time \(t\), extend the holes by making all \(\ell\) slots outside
\(\Lambda\) available.  Let \(\mathcal H_t\) be the resulting infinite
available-slot graph, with all pairs of available slots above each
interaction bond joined.  For a finitely supported real function \(v\)
on this graph, define
\begin{equation}\label{eq:full-energy}
 \mathcal D_t(v)=
 \sum_{e=\{i,j\}\in E(\mathcal H_t)}c_e\bigl(v(i)-v(j)\bigr)^2,
 \qquad c_e=\tfrac12 J(x_j-x_i).
\end{equation}
Each unoriented edge is counted once.  A killed vector is extended by
zero on killed slots, including the boundary and the exterior.
On a smooth interval, in elapsed time \(u\) in either traversal direction,
\begin{equation}\label{eq:norm-dissipation}
 \frac{d}{du}\|v(u)\|_2^2=-2\mathcal D_{t(u)}(v(u)).
\end{equation}
To verify this, pair the symmetric jump generator with \(v\); each
edge contributes \(-2c_e(v(i)-v(j))^2\).  Edges to killed slots retain
their exit rates and give the same expression with the killed value
zero.  There is no bond from \(\Lambda^\circ\) directly to the exterior,
so the infinite extension creates no extra term.  The same identity
holds after additional killing outside a fixed spatial set.  Squared
norm cannot increase at any instantaneous factor.

\subsection{Regular anchors and local diffusion}

For a coset \(\gamma\), let \(\mathcal H_t^\gamma\) contain its
available slots and only the nearest-neighbor edges in bar coordinates.
Write
\[
 \mathcal D_t^\gamma(v)
   =\frac12\sum_{\{i,j\}\in E(\mathcal H_t^\gamma)}
                   (v(i)-v(j))^2,
 \qquad c_0=\min_{1\le j\le3}J(a_j)>0.
\]
The graphs for different cosets have disjoint vertex sets, and their
edges are a subset of the full interaction edges.  Therefore
\begin{equation}\label{eq:energy-comparison}
 \mathcal D_t(v)\ge c_0
       \sum_{\gamma\in\mathfrak C}\mathcal D_t^\gamma(v).
\end{equation}
This inequality is the only place where the geometric subgraphs replace
the full graph; in particular it changes constants, not the diffusion
matrix that will determine Bloch's coefficient.

Fix
\[
 \alpha=\frac1{100},\qquad \nu=\frac1{300},\qquad
 L\ge\max\{6,\ell|\mathfrak C|\}.
\]
\begin{definition}[Sparse collections]\label{def:sparse}
A finite collection of slot points with site coordinates \((x_i)\) is
\emph{sparse} if
\begin{equation}\label{eq:sparse-growth}
 \#\{i:\bar x_i\in B_h(z)\}\le L(1+h)^\alpha
 \qquad(z\in\Z^3,\ h\ge0).
\end{equation}
Repeated bar coordinates count with their multiplicities, including
points from different cosets.  The same definition applies to a list of
centers in \(\Z^3\).  Every subcollection remains sparse.
\end{definition}

In coset \(\gamma\), call a coordinate site \emph{blocked} if it has
no available slot; denote this finite set by \(\mathcal O_t^\gamma\).
Join blocked sites whose supremum distance is at most eight.
For \(w\in\mathcal O_t^\gamma\), let \(\mathcal C_t^\gamma(w)\)
be its component and put
\[
 \operatorname{rad}_w\mathcal C_t^\gamma(w)
   =\max_{v\in\mathcal C_t^\gamma(w)}|v-w|_\infty.
\]
Boundary sites and all sites outside \(\Lambda\) have every slot
available.  Partial exposure at a site is permitted: an unblocked site
can have any number of available slots between one and \(\ell\).

We recall the precise local-connectivity input from
\cite[Lemma~5.5]{Ordering}.  For every fixed integer \(b\ge1\), there
is an integer \(\rho(3,b)\ge b+1\) such that if
\(u,v\notin F\subset\Z^3\), \(|u-v|_\infty\le b\), and
\(\#(F\cap B_{\rho(3,b)}(u))<6\), then a nearest-neighbor path from
\(u\) to \(v\) avoids \(F\), stays in that ball, and has length at
most \((2\rho(3,b)+1)^3\).  Fix an integer, as in the radius choice of
\cite[Equation~(5.13)]{Ordering},
\begin{equation}\label{eq:geometric-radius}
 M_*\ge\max\{48,\rho(3,44)+2\}.
\end{equation}
The numbers 44 and 48 cover the bounded connections around components
of at most five blocked sites.  Fix a lower threshold \(R_0\) large
enough for all subsequent deterministic estimates; it includes
\(R_0\ge M_*+47\).  This threshold depends only on the fixed geometry.

\begin{definition}[Regular slice]\label{def:regular}
A slice \(t\) is \emph{regular for the anchor \(z\) at radius \(R\)}
if the following hold in every coset \(\gamma\):
\begin{enumerate}
 \item for every \(w\in B_{11R}(z)\),
 \(\#(\mathcal O_t^\gamma\cap B_{M_*}(w))<6\);
 \item for every \(w\in\mathcal O_t^\gamma\setminus B_{10R}(z)\),
 \(\operatorname{rad}_w\mathcal C_t^\gamma(w)
       \le(1+|w-z|_\infty)^\nu\).
\end{enumerate}
The anchor itself need not be an available site.
\end{definition}

These are precisely the regular-anchor conditions in
\cite[Definition~5.3]{Ordering}, on each coordinate lattice.  They are
conditions on a single deterministic graph; no pin probability estimate
is assumed in the next three propositions.

\begin{proposition}[Local Nash inequality]\label{prop:nash}
Suppose \(R\ge R_0\) and \(t\) is regular for \(z\) at radius \(R\).
Every nonnegative function \(v\) on \(\mathcal H_t\), supported on
slots above \(\mathfrak B_R(z)\), satisfies
\begin{equation}\label{eq:local-nash}
 \|v\|_2^{10/3}\le C\mathcal D_t(v)\|v\|_1^{4/3}.
\end{equation}
The constant is independent of \(R,z,t\) and the unavailable slots.
\end{proposition}

\begin{proof}
We give the transfer argument of \cite[Proposition~5.6]{Ordering},
followed by the lattice averaging proof of Nash's inequality; the
energy-to-diffusion method originates in \cite{Nash1958}.
First work in one coset and write its bar coordinates simply as sites.
At each unblocked site \(w\), choose a slot where \(v\) is maximal
and call this value \(V(w)\).  Each blocked site of \(B_{R+3}(z)\)
has a surviving nearest neighbor: otherwise its six neighbors would
contradict local regularity.  Choose one, denoted \(\phi(w)\), and define
\[
 f(w)=
 \begin{cases}
 V(w),&w\notin\mathcal O_t^\gamma,\\
 V(\phi(w)),&w\in\mathcal O_t^\gamma\cap B_{R+3}(z),\\
 0,&w\in\mathcal O_t^\gamma\setminus B_{R+3}(z).
 \end{cases}
\]
Then \(f\) is supported on \(B_{R+1}(z)\), and
\begin{equation}\label{eq:nash-norm-transfer}
 \|v\|_2^2\le\ell\|f\|_2^2,
 \qquad \|f\|_1\le7\|v\|_1.
\end{equation}
For a lattice edge with a nonzero endpoint value of \(f\), both
values come from maximizing slots at surviving sites at distance at
most three from each other, within \(B_{R+3}(z)\).
The local-connectivity input with \(b=44\) joins these sites by a path of fixed length
and displacement avoiding blocked sites.  Lift it to their maximizing
slots.  Telescoping along the lifted path and using Cauchy--Schwarz
bounds the squared difference across the original edge by a fixed
multiple of the sum of squared slot differences on that path.
Any slot edge can occur in only a bounded number of these paths,
since their original edges lie in a fixed neighborhood of it.  Hence
\begin{equation}\label{eq:nash-energy-transfer}
 \mathcal E_{\Z^3}(f):=
 \sum_{\{w,w'\}:|w-w'|_1=1}(f(w)-f(w'))^2
 \le C\mathcal D_t^\gamma(v).
\end{equation}

For completeness, average \(f\) over its translates by
\(\{0,\ldots,b-1\}^3\), and denote the result by \(A_bf\).
Telescoping coordinate translations and using the triangle inequality
in \(\ell^2\) gives
\(\|f-A_bf\|_2^2\le Cb^2\mathcal E_{\Z^3}(f)\).
The averaging kernel has \(2\)-norm \(b^{-3/2}\), so
\(\|A_bf\|_2^2\le b^{-3}\|f\|_1^2\).  Therefore
\[
 \|f\|_2^2\le Cb^2\mathcal E_{\Z^3}(f)+2b^{-3}\|f\|_1^2.
\]
For \(f\ne0\), take
\(b=\lceil(4\|f\|_1^2/\|f\|_2^2)^{1/3}\rceil\).
The second term is at most half the left side; since
\(\|f\|_1\ge\|f\|_2\), rounding changes only the constant.
Absorbing this term proves
\(\|f\|_2^{10/3}\le C\mathcal E_{\Z^3}(f)\|f\|_1^{4/3}\).
The zero case is immediate.  Equations
\eqref{eq:nash-norm-transfer}--\eqref{eq:nash-energy-transfer} give
\eqref{eq:local-nash} for one coset with \(\mathcal D_t^\gamma\).
Finally, writing \(v_\gamma\) for the restriction to coset \(\gamma\), use
\[
 \Bigl(\sum_\gamma\|v_\gamma\|_2^2\Bigr)^{5/3}
 \le |\mathfrak C|^{2/3}\sum_\gamma\|v_\gamma\|_2^{10/3},
 \qquad \|v_\gamma\|_1\le\|v\|_1,
\]
and \eqref{eq:energy-comparison} to obtain the assertion on all cosets.
\end{proof}

\subsection{Flows and sparse sources}

The Nash inequality controls a mass while it evolves.  We next control
its pairing with an arbitrary function at a single slice.  A flow with
divergence equal to that mass gives precisely this bound by summation
by parts.  Fix an orientation \((e^-,e^+)\) for every unoriented slot
edge, and define divergence as outgoing minus incoming flow.  If \(F\)
is a flow and \(Y\) has finite support, then
\begin{equation}\label{eq:flow-pairing}
 \langle\Div F,Y\rangle
 =\sum_eF(e)\bigl(Y(e^-)-Y(e^+)\bigr).
\end{equation}
For an edge in one coset and a coordinate site \(y\), write
\(\dist(e,y)\) for the minimum supremum distance from \(y\) to an
endpoint coordinate of \(e\).

We state the deterministic input, including the features of its
construction needed for a sharper estimate below.  Its hypotheses do
not refer to an exposure law or to a relation between \(R\) and a time
scale.

\begin{lemma}[Routed unit flows: geometric input]\label{lem:routed-flow}
Consider an available-slot nearest-neighbor graph on \(\Z^3\), with
at most \(\ell\) slots per site, all pairs of available slots above
each nearest-neighbor site edge joined, and all slots available outside
a finite set.  Suppose \(z\) satisfies the two conditions in
Definition~\ref{def:regular} for this graph, and \(R\ge R_0\).
For each available slot \(a\) at \(y\in B_R(z)\), there is a flow
\(F_a\) such that
\begin{equation}\label{eq:global-flow-envelope}
 \Div F_a=\ind{\{a\}},\qquad
 |F_a(e)|\le C(1+\dist(e,y))^{-k},\qquad
 k=2-3\nu=\frac{199}{100}.
\end{equation}
The flows can be chosen with the following common routing construction.
For each blocked component \(C\), thicken it by supremum radius four,
fill the finite nearest-neighbor components of the complement, and call
the resulting finite hull \(H_C\).  Let \(\operatorname{Box}_C^+\)
be the unit enlargement of the coordinate bounding box of that
thickening.  The hull lies in this box.  Keep the inclusion-maximal
hulls, which are disjoint, and send each original site in such a hull
to a chosen surviving exterior-boundary representative, fixing other
sites.  Apply this map only once to each original endpoint.

There are simple surviving routes for the original lattice edges.
Each route is contained in its original edge together with the boxes
\(\operatorname{Box}_C^+\) of the maximal hulls containing its original
endpoints.  Lift the routes using a fixed available slot at every
surviving site.  The flow \(F_a\) is the sum of a unit path flow along a connector
of uniformly bounded length and displacement from \(a\), and the
routed values of a lattice
unit flow \(G_y\) satisfying
\begin{equation}\label{eq:lattice-flow-envelope}
 |G_y(e')|\le C(1+\dist(e',y))^{-2}.
\end{equation}
All constants are uniform over the graph, anchor, radius, and source.
\end{lemma}

This is \cite[Proposition~6.2]{Ordering} together with its stated
construction: hull containment and surviving boundary routes are in
\cite[Lemma~6.1 and Section~6.1]{Ordering}; the once-applied endpoint
map and the routes, including adjacent maximal hulls, follow
Equations~(6.2)--(6.3) there; the bounded connector, lattice flow, and
lifting are in the proof of Proposition~6.2, Equations~(6.6)--(6.8).
These statements allow arbitrary available-slot counts between one and
\(\ell\) at unblocked sites.  Thus they apply on every coset here,
whose unavailable slots form a finite set.  No assertion about the
nearest-neighbor spin model is being used.

Two such flows, whose source sites are \(y,y'\), satisfy
\begin{equation}\label{eq:flow-overlap}
 \sum_e|F_a(e)F_b(e)|
 \le C(1+|y-y'|_\infty)^{-q},
 \qquad q=2k-3=\frac{49}{50}.
\end{equation}
Here the flows may use different regular anchors.  To check this estimate,
put \(H=1+|y-y'|_\infty\).  In a ball of radius
\(|y-y'|_\infty/3\) about either source, one envelope is at most
\(CH^{-k}\), and the sum of the other is at most \(CH^{3-k}\).
Away from those two balls, dyadic shells of radii at least a fixed
multiple of \(H\) give \(C\sum_{j\ge0}(2^jH)^{3-2k}\).
Both bounds are \(CH^{3-2k}\), since \(3/2<k<3\).
Bounded \(H\) follows by square summability.  There are only a bounded
number of slot edges per coordinate site, so the same calculation
applies to edge sums.  This also proves the overlap estimate recorded
in \cite[Lemma~6.3]{Ordering}.

\begin{proposition}[Energy of sparse sources]\label{prop:sparse-energy}
Let \(R\ge R_0\), and let \((z_i)_{i=1}^m\) be a sparse list of
centers, counting repetitions.  At a fixed slice let \(I\subseteq
\{1,\ldots,m\}\) be any subset for which the slice is regular for
\(z_i\) at radius \(R\).  For \(i\in I\), let \(\mu_i\ge0\)
be a mass on the available slots, supported above \(\mathfrak B_R(z_i)\)
and with total mass at most one.  For every finitely supported real \(Y\),
\begin{equation}\label{eq:sparse-source-pairing}
 \left|\left\langle\sum_{i\in I}\mu_i,Y\right\rangle\right|^2
 \le C m R^\alpha\mathcal D_t(Y).
\end{equation}
\end{proposition}

\begin{proof}
In each coset sum the unit flows of Lemma~\ref{lem:routed-flow} with
the coefficients supplied by the restrictions of the \(\mu_i\),
and call the result \(F^\gamma\).  Its divergence is the sum of
these restrictions.  The supports of the masses are finite.
If \(|z_i-z_j|_\infty>4R\), the two source sites in their respective
supports have distance at least \(|z_i-z_j|_\infty/2\); otherwise
use the uniform constant overlap bound.  Equation~\eqref{eq:flow-overlap}
and the mass bound at each center give
\[
 \sum_e|F^\gamma(e)|^2
 \le C\sum_{i\in I}\left[
 \#\{j\in I:|z_j-z_i|_\infty\le4R\}
 +\sum_{\substack{j\in I\\|z_j-z_i|_\infty>4R}}
          (1+|z_j-z_i|_\infty)^{-q}\right].
\]
The first count is at most \(CL R^\alpha\).  Divide the second
sum into annuli \(2^n4R<|z_i-z_j|_\infty\le2^{n+1}4R\);
its bound is
\(CL\sum_{n\ge0}(2^nR)^{\alpha-q}\le C\), because \(q>\alpha\)
and \(R\ge1\).  Consequently
\(\sum_\gamma\sum_e|F^\gamma(e)|^2\le CmR^\alpha\).
Using \eqref{eq:flow-pairing}, Cauchy--Schwarz on the disjoint coset
edge sets, and \eqref{eq:energy-comparison}, we get
\[
 \left|\left\langle\sum_{i\in I}\mu_i,Y\right\rangle\right|^2
 \le CmR^\alpha\sum_\gamma2\mathcal D_t^\gamma(Y)
 \le CmR^\alpha\mathcal D_t(Y).
\]
This includes the empty set \(I\), for which the flow is zero.
\end{proof}

\subsection{The sharper bound for a spread source}

The next estimate converts a \(2\)-norm bound \(C_1\tau^{-3/4}\)
on a subunit mass into a flow-energy bound \(C\tau^{-1/2}\).
The global exponent \(k<2\) in
\eqref{eq:global-flow-envelope} alone loses a power of \(\tau\).
Local regularity repairs this loss: near the source, every obstacle
that the routing encounters has bounded size.

\begin{lemma}[Decay up to the localization radius]\label{lem:near-flow}
The same flows as in Lemma~\ref{lem:routed-flow} satisfy
\begin{equation}\label{eq:near-flow-envelope}
 |F_a(e)|\le C(1+h)^{-2}
 \qquad\text{if }h=\dist(e,y)\le R.
\end{equation}
The constant and the lower threshold \(R_0\) are independent of \(R\).
\end{lemma}

\begin{proof}
A blocked component meeting \(B_{10R}(z)\) has at most five sites.
Indeed, an exploration of six sites starting at a point in that ball
would place them all in a radius-40 ball, contradicting regularity and
\(M_*\ge48\).  Thus such a component has diameter at most 32.
Its radius-four thickening has bounding-box diameter at most 40, so
\(\operatorname{Box}_C^+\) has diameter at most 42.

If a component avoids \(B_{10R}(z)\), put
\(b_C=\min_{w\in C}|w-z|_\infty>10R\).
Remote regularity, applied at a minimizer, gives
\[
 \operatorname{diam}_\infty(\operatorname{Box}_C^+)
 \le 2(1+b_C)^\nu+10\le b_C/10
\]
after increasing the fixed threshold \(R_0\).
The minimizing site belongs to the box, so the entire box lies outside
\(B_{9R}(z)\).

Now let \(e\) be a target slot edge with \(h\le R\).  Both its
endpoint sites lie in \(B_{2R+1}(z)\).  If a routed original edge
\(e'\) uses \(e\), route containment in Lemma~\ref{lem:routed-flow}
says that this occurrence is either on \(e'\) itself or in an
incident enlarged box.  That box cannot be remote, and therefore has
diameter at most 42.  Since it contains an original endpoint of \(e'\),
the edge \(e'\) is within a fixed distance of \(e\).  Thus only a
bounded number of original edges contribute to \(e\), and their
source distances differ from \(h\) by at most a fixed constant.
Each route is simple and uses \(e\) at most once.  Summing
\eqref{eq:lattice-flow-envelope} over these edges gives
\eqref{eq:near-flow-envelope}; the bounded connector changes only the
constant.  This argument also covers routes between adjacent maximal
hulls, since their middle portion is the original edge with reversed
orientation.
\end{proof}

\begin{proposition}[Spread-source energy]\label{prop:spread-energy}
Suppose \(R\ge R_0\), \(1\le\tau\le R\), and the slice \(t\)
is regular for \(z\) at radius \(R\).
Let \(\mu\ge0\) be supported on available slots above
\(\mathfrak B_R(z)\), with
\[
 \|\mu\|_1\le1,\qquad \|\mu\|_2\le C_1\tau^{-3/4}.
\]
For every finitely supported real function \(Y\) on \(\mathcal H_t\),
\begin{equation}\label{eq:spread-source-pairing}
 |\langle\mu,Y\rangle|^2
 \le C\tau^{-1/2}\mathcal D_t(Y).
\end{equation}
The constant may depend on \(C_1\) and the fixed geometry, but is
independent of \(\tau,R\), the slice, and the volume.
\end{proposition}

\begin{proof}
Work first on one coset, with the restriction of \(\mu\), and set
\(F=\sum_a\mu(a)F_a\).  Then \(\Div F=\mu\).
For each source at \(y_a\), split its flow at distance
\(s=\sqrt\tau\) into near and far parts.  The parts need not have
prescribed divergences; their sum is the original flow.

Index each edge by one endpoint site, with bounded multiplicity, and
put \(\bar\mu(w)=\sum_{a:y_a=w}\mu(a)\).  Since at most \(\ell\)
slots lie over a site,
\(\|\bar\mu\|_2\le\sqrt\ell\|\mu\|_2\).
As \(s\le R\), Lemma~\ref{lem:near-flow} bounds the near part
pointwise by a convolution with a nonnegative lattice kernel
\[
 \kappa_s(w)=C(1+|w|_\infty)^{-2}
                     \ind{\{|w|_\infty\le s+1\}},
 \qquad \|\kappa_s\|_1\le Cs.
\]
Here changing the minimum distance to an edge into the distance from
its indexed endpoint costs at most one.  Young's inequality yields
\begin{equation}\label{eq:spread-near-energy}
 \|F_{\mathrm{near}}\|_2
 \le C\|\kappa_s*\bar\mu\|_2
 \le Cs\|\mu\|_2
 \le C\tau^{-1/4}.
\end{equation}
For the far part of each individual flow, count at most \(C(1+h)^2\)
slot edges at distance \(h\), use the exponent 2 until \(R\), and
then use the exponent \(k\).  This gives
\begin{align}
 \sum_{e:\,\dist(e,y_a)>s}|F_a(e)|^2
 &\le C\sum_{s<h\le R}(1+h)^{-2}
      +C\sum_{h>R}(1+h)^{2-2k} \notag\\
 &\le C\bigl(\tau^{-1/2}+R^{-49/50}\bigr)
 \le C\tau^{-1/2}.                         \label{eq:spread-far-energy}
\end{align}
The last comparison uses exactly \(R\ge\tau\ge1\) and
\(49/50>1/2\).  The triangle inequality in edge \(\ell^2\), with
\(\sum_a\mu(a)\le1\), bounds
\(\|F_{\mathrm{far}}\|_2\le C\tau^{-1/4}\).
Together with \eqref{eq:spread-near-energy}, this proves
\(\sum_e|F(e)|^2\le C\tau^{-1/2}\).

Apply this construction on every coset.  Each restriction satisfies
the same mass and \(2\)-norm bounds, so the sum of all coset flow
energies is at most \(C|\mathfrak C|\tau^{-1/2}\).
Summation by parts \eqref{eq:flow-pairing}, Cauchy--Schwarz, and
\eqref{eq:energy-comparison} now prove
\eqref{eq:spread-source-pairing}.
\end{proof}

In particular, when \(\tau=K\beta\) and
\(R=\lceil a(K+1)\beta\rceil\) with fixed \(a\ge1\), the constant
in Proposition~\ref{prop:spread-energy} does not depend on \(K\).
All statements of this section are deterministic at a single exposure
or slice.  Sources and retained subsets may therefore depend on the
entire exposure.  The estimates require no measurable selection of
flows: the conclusions are inequalities between the given masses and
energies.
